% Part: computability
% Chapter: computability-theory
% Section: def-functions-self-reference

\documentclass[../../../include/open-logic-section]{subfiles}

\begin{document}

\olfileid{cmp}{thy}{slf}
\olsection{د ځان-مراجعې په کارولو د تابعو تعريف}

په عمومي ډول دا ګټوره ده چې تابعې د هغو په خپلو اصطلاحاتو کښې تعريف
کړے شو۔ د بېلګې په توګه، ورکړو محاسبه کېدونکو تابعو~$k$، $l$ او~$m$
ته د ثابت ټکي لمه وايي چې داسې جزوي محاسبه کېدونکې تابع~$f$ شته چې
د هر~$y$ دپاره دا معادله پوره کوي:
\[
f(y) \simeq
\begin{cases}
k(y) & \text{که $l(y) = 0$ وي} \\
f(m(y)) & \text{که داسې نۀ وي۔}
\end{cases}
\]
په لا ځانګړي ډول، ~$f$ بيا داسې تر لاسه کېږي:
\[
g(x,y) \simeq
\begin{cases}
k(y) & \text{که $l(y) = 0$ وي} \\
\cfind{x}(m(y)) & \text{که داسې نۀ وي}
\end{cases}
\]
او بيا د ثابت ټکي لمه کارول کېږي څو داسې شاخص~$e$ وموندل شي چې
$\cfind{e}(y) \simeq g(e,y)$ وي۔
\textbf{د ژباړې د نسخې سمون (OLCMP-044):} سرچينې د ثابت ټکي په دې
پليينه کښې د دوو جزوي تابعو د ممکنه ناتعريفو ارزښتونو تر منځ عادي
برابري کارولې وه؛ دلته د قضيې ټاکلې جزوي برابري ساتل شوې ده۔

د يوې محسوسې بېلګې دپاره د ``تر ټولو لوی ګډ مقسوم عليه'' تابع
$\fn{gcd}(u,v)$ داسې تعريفېدے شي:
\[
\fn{gcd}(u,v) \simeq
\begin{cases}
v & \text{که $u = 0$ وي} \\
\fn{gcd}(\fn{mod}(v, u), u) & \text{که داسې نۀ وي}
\end{cases}
\]
چېرته چې~$\fn{mod}(v, u)$ پر~$u$ د~$v$ د وېش پاتې شونې څرګندوي۔ د
ثابت ټکي لمې ته مراجعه ښيي چې~$\fn{gcd}$ جزوي محاسبه کېدونکې ده۔
(په حقيقت کښې دا په پورته بڼه راوستل کېدے شي، چې~$y$ د جوړې
$\tuple{u, v}$ کوډ وګڼل شي۔) بيا پر~$u$ استقرا ښيي چې په حقيقت کښې
~$\fn{gcd}$ هرځاے تعريف شوې ده۔

البته داسې ځان-مراجع تعريفونه جوړولے شو چې له پورته بېلګو ډېر پېچلي
وي۔ زياتره پروګرامي ژبې په يوه يا بله طريقه د تابعو تعريف په خپلو
اصطلاحاتو کښې مني۔ پام وکړئ چې دا د هغې وړتيا په پرتله لږ حيرانوونکے
دے چې يوه تابع د هغه الګوريتم د \emph{شاخص} په وسيله تعريف کړو چې
هماغه تابع محاسبه کوي؛ او د ثابت ټکي قضيه په بشپړ عموميت کښې همدا
اجازه راکوي۔
\textbf{د ژباړې د نسخې سمون (OLCMP-037):} سرچينې د دې وروستۍ جملې
په پای کښې د يو الګوريتم له خوا محاسبه کېدونکې يوه تابع په جمع نوم
ياده کړې وه۔ دلته شمېر د هماغې ټاکلې تابع سره يو شان ساتل شوے دے۔

\end{document}
